\chapter{Introduction}
\label{ch:introduction}

% TODO: Open the thesis with a hook that motivates dynamics-constrained
% manipulation and previews the arc from sampling-based to generative methods.

\section{The State of Robot Manipulation Today}
\label{sec:intro:state}

% TODO: Frame the contemporary landscape of robot manipulation.

\subsection{Progress in Pipelined Industrial Systems}
\label{sec:intro:industrial}

% TODO: Discuss reliability and debuggability as the dominant priorities of
% deployed industrial manipulation pipelines.

\subsection{Progress in Academic Systems}
\label{sec:intro:academic}

% TODO: Discuss academic systems pushing toward generalist manipulation
% behaviors.

\subsection{Common Limitations}
\label{sec:intro:limitations}

% TODO: Brittleness, task-specific engineering, and poor cross-task
% generalization remain the shared bottlenecks across both styles of system.

\section{The Geometric vs.\ Dynamic Divide}
\label{sec:intro:divide}

% TODO: Most planning reasons over collisions, grasps, and singularities ---
% all purely geometric quantities.

% TODO: Real-world tasks demand reasoning over contact accelerations,
% traversal time, and impulsive forces.

% TODO: Argue that accessing the full range of manipulation requires
% integrating dynamics constraints into motion generation.

\section{Thesis Statement and Contributions}
\label{sec:intro:contributions}

% TODO: State the thesis: a progression of methods --- sampling-based to
% generative --- for dynamics-constrained motion generation.

% TODO: Enumerate contributions across tabletop manipulation, actuation
% failures, liquid handling, and payload transport.

\section{Roadmap of the Thesis}
\label{sec:intro:roadmap}

% TODO: Walk the reader through the structure of the remaining chapters.
